
$\varepsilon$ l'homomorphisme d'augmentation : $K^{\bullet}(X)\longrightarrow H^0(X,\underline{\mathbb Z})$.

\textbf{Proposition 6.1.} \uline{Soit} $X$ \uline{un schéma quasi}-\uline{compact ayant un faisceau ample} $\underline{L}$. \uline{Alors} $\Fil^1K^{\bullet}(X)$ \uline{est un nilidéal}.

Comme $\Fil^1K^{\bullet}(X)$ est engendré par les éléments du type $E-\varepsilon(E)$, où $E$ est la classe d'un $\underline{O}_X$-Module localement libre de type fini $\underline{E}$, et $\varepsilon:K^{\bullet}(X)\to H^0(X,\underline{\mathbb Z})$ l'augmentation de $K^{\bullet}(X)$, il suffit de montrer qu'un tel élément est nilpotent. La quasi-compacité de $X$, et la commutation du $K^{\bullet}$ aux sommes directes finies permettent de se ramener au cas où $\underline{E}$ est de rang constant. Enfin, en passant au fibré de drapeaux de $\underline{E}$, ce qui donne une injection sur les $K^{\bullet}$ (3.1), et conserve l'hypothèse d'existence d'un faisceau ample, on peut supposer $\underline{E}$ inversible.

Soit $\underline{E}^*$ le dual de $\underline{E}$. Pour $n$ assez grand, $\underline{E}^*(n)=\underline{E}^*\otimes_{\underline{O}_X}\underline{L}^{\otimes n}$ est engendré par ses sections globales; on peut donc trouver un entier $m$, et un homomorphisme surjectif : $(\underline{O}_X)^m\longrightarrow \underline{E}^*(n)$, donc (tensorisant par $\underline{E}(-n)$) un homomorphisme surjectif : $(\underline{E}(-n))^m\longrightarrow\underline{O}_X$. Le noyau de cet homomorphisme est localement libre de rang $m-1$. On obtient donc, en notant $L$ la classe de $\underline{L}$:
\[
  \lambda^m(mEL^{-n}-1)=0 .
\]
Or $\lambda^m(mEL^{-n}-1)=(-1)^m\lambda_{-1}(mEL^{-n})=(-1)^m\lambda_{-1}(EL^{-n})^m=(-1)^m(1-EL^{-n})^m$. Par conséquent, pour tout faisceau inversible $\underline{E}$, l'élément $1-EL^{-n}$ est nilpotent pour $n$ assez grand, ou encore son produit $L^n-E$ par $L^n$ est nilpotent.

On peut alors écrire:
\[
  E-1=(L^n-1)-(L^n-E)
\]
et d'après ce qui précède, $L^n-E$ et $L^n-1$ sont nilpotents si on choisit $n$ assez grand. Donc $E-1$ est nilpotent.


